\documentclass[12pt]{article}
\usepackage[utf8]{inputenc}
\usepackage[T1]{fontenc}
\usepackage{hyperref}
\title{Wintu Language Resources for K-12 Curriculum Development}
\author{Compiled by Claude Sonnet 5.5 (claude-sonnet-5-5) for Andr\'as Kornai}
\date{}
\begin{document}
\maketitle

\section{Introduction}
Wintu (ISO 639-3: wit; Glottocode wint1259) is a Wintuan language of northern California. According to Wikipedia it was considered extinct after the death of Flora Jones in 2003, and one partial speaker was reported in 2011. A revitalization effort is under way within the Winnemem Wintu. This report is based on one web search. Only resources that appeared in the results are listed, and little concrete material was found.

\section{Primary Educational Resources}
No published primer, textbook or school programme appeared in the results. Omniglot reports that the Winnemem Wintu Language Project has worked since 2011 with linguist Stefan Liedtke on an alphabet, learning materials and a talking dictionary. I found no confirmation that the dictionary was completed or released.

\section{Linguistic and Reference Materials}
\begin{itemize}
\item Wikipedia article on the Wintu language. \url{https://en.wikipedia.org/wiki/Wintu_language}
\item Omniglot page on Wintu, including its writing system. \url{https://omniglot.com/writing/wintu.htm}
\end{itemize}

\section{Religious and Community Resources}
\begin{itemize}
\item Wikipedia article on the Winnemem Wintu. It names Headman Marc Franco as working with the Indigenous Language Institute on the language as of 2011. \url{https://en.wikipedia.org/wiki/Winnemem_Wintu}
\end{itemize}

\section{Digital and Audio Resources}
The Endangered Languages Project lists resources for Wintu, including a Winnemem-language video story dated 28 August 2025. \url{https://endangeredlanguages.com/node/37667/resources}

\section{Curriculum Development Context}
The language has very few speakers and the revitalization effort is recent. A curriculum should be developed with the Winnemem Wintu Tribe and its language project. The tribe's own website was not seen in the results and would need to be located. Archival linguistic documentation was not identified in this search and should be sought separately.

\end{document}